\chapter{Open Problems in Ruliometric Science}
\label{sec:open-problems-ruliometric}

This is the part where the book stops explaining and starts throwing you work.

We’ve built:

\begin{itemize}
  \item an ontology (RMG + DPO as substrate),
  \item a geometry (MRMW, worldlines, curvature),
  \item a physics (CΩSMOS as an RMG universe),
  \item machines (time-travel debuggers, counterfactual engines, MORIARTY),
  \item an architecture (CΩMPILER, runtime, storage),
  \item and a sketch of intelligence + ethics in this setting.
\end{itemize}

That’s not a finished theory. It’s a coordinate system on a space that barely has a name.

Let’s name it now:

\begin{quote}
\textbf{Ruliometric science} is the quantitative study of MRMW---the geometry of the space of all possible computations, rule systems, and universes---and the machines that move through it.
\end{quote}

Think: differential geometry + complexity theory + dynamical systems + statistical physics, all stapled to a combinatorial substrate and pointed at everything from physics to AI to institutions.

This chapter is a non-exhaustive hit list of open problems in that field.
Some are crisp conjectures, some are vague research programs. All of them are
deliberately pointed: things where progress would change how we build systems.

You can read this as a roadmap, a shopping list, or an invitation.

\subsubsection*{30.1  Make MRMW a Proper Geometric Object}

Right now MRMW---the space of all RMGs + rule sets + histories---is described
hand-wavily as a "phase space of all computations."

That’s not good enough. We need actual math.

\paragraph{Problem 1: Define usable metrics on MRMW.}

We want distances:

\begin{itemize}
  \item between two states $(G_1,\mathcal{R}_1)$ and $(G_2,\mathcal{R}_2)$;
  \item between two rule sets $\mathcal{R}_1, \mathcal{R}_2$ (rulial distance);
  \item between two worldlines (histories) $\gamma_1, \gamma_2$.
\end{itemize}

Constraints:

\begin{itemize}
  \item invariant under obvious symmetries (relabelings, coarse-grainings),
  \item computable / approximable for realistic systems,
  \item meaningful: nearby points correspond to "similar behavior" in some observable sense.
\end{itemize}

Concrete questions:

\begin{enumerate}
  \item Can we define a Gromov–Hausdorff-type distance between causal RMG universes?
  \item Can we define a Wasserstein-like distance between distributions over worldlines?
  \item Can we parameterize "distance in rule space" by how much an agent’s value functional changes under compilation of the same spec?
\end{enumerate}

\paragraph{Problem 2: Curvature in MRMW that actually does something.}

We’ve spoken about curvature as "hardness" or "bending’’ of computational paths.
Turn that into real invariants:

\begin{itemize}
  \item define curvature tensors or scalar curvatures on local neighborhoods in MRMW;
  \item relate curvature to:
    \begin{itemize}
      \item algorithmic complexity (sharp curvature $\Leftrightarrow$ NP-ish behavior?),
      \item sensitivity to initial conditions (chaos),
      \item robustness of coarse-grainings.
    \end{itemize}
\end{itemize}

Yes, this sounds like re-inventing Ricci flow on computation space. Good. Do it.

\paragraph{Problem 3: Topology of MRMW.}

What are the large-scale features?

\begin{itemize}
  \item Are there connected components of rule sets that cannot be deformed into one another without crossing phase transitions (e.g., from reversible to irreversible dynamics)?
  \item Are there "holes" corresponding to classes of universes that can’t be approximated by finite RMGs or local DPO rules?
  \item Are there universal attractors: regions most random rule sets flow toward under some natural renormalization?
\end{itemize}

If you’re a topologist who’s bored of manifolds, here’s your playground.

\subsubsection*{30.2  Complexity as Geometry: Hard Problems, Hard Curvature}

We pretend complexity theory is about Turing machines and oracle models.
In Ruliometry, it should be about geometry.

\paragraph{Problem 4: Geometric characterization of complexity classes.}

Conjecture-level idea:

\begin{quote}
Computational complexity classes correspond to families of regions in MRMW with characteristic curvature, volume-growth, and geodesic-length profiles.
\end{quote}

That is:

\begin{itemize}
  \item P-like behavior: geodesics of length polynomial in input size between typical start and goal states.
  \item NP-like behavior: exponential explosion of geodesic count, but existence of at least one short path.
  \item BQP-like behavior: cheap access to whole bundles of paths via quantum-like superposition rules.
\end{itemize}

Questions:

\begin{enumerate}
  \item Can we define an invariant $K$ on problem encodings in MRMW that lower-bounds time/space complexity?
  \item Can we prove analogues of known complexity separations or collapses as curvature / volume constraints?
  \item Is there a notion of average-case curvature that correlates with phase transitions in solvability?
\end{enumerate}

If any of this works, "complexity theory" becomes a chapter in Ruliometry, not a separate cult.

\paragraph{Problem 5: Local NP collapse and curvature.}

We claimed in Part~III that certain local regions of MRMW can exhibit effective "NP collapse" due to geometric properties (e.g., constrained search regions, strong pruning by rewrite laws).

Formalize that as:

\begin{quote}
Given a subregion $\Omega \subset \text{MRMW}$ and a problem encoding $\Pi$ supported in $\Omega$, define conditions on curvature/volume of $\Omega$ that imply tractable solution of $\Pi$.
\end{quote}

If you can prove non-trivial examples, you get:

- structural explanations for why certain practically-hard problems are easy in real-world instances,
- design rules for building machines that sit in those nice regions.

\subsubsection*{30.3  Classification of Universes: A Zoo of RMGs}

Right now, we have a few toy universes (cellular automata, graph rewrites) and one big scary universe (ours).

We need a \emph{taxonomy}.

\paragraph{Problem 6: Universality classes of RMG rule sets.}

We want to cluster rule sets $\mathcal{R}$ into universality classes, not just by Turing-completeness, but by:

\begin{itemize}
  \item emergent dimension,
  \item symmetry structure,
  \item typical complexity scaling,
  \item stability under coarse-graining.
\end{itemize}

Questions:

\begin{enumerate}
  \item Can we define RG-like flows on rule space that coarse-grain a rule set into a simpler effective theory?
  \item Do different microscopic rules flow to the same macroscopic universality class?
  \item Is there a "periodic table" of universe types under these flows?
\end{enumerate}

\paragraph{Problem 7: Dualities between universes.}

Two very different rule sets can yield the same macroscopic physics (or the same computation) under different coarse-grainings.

\begin{quote}
Formalize when $(G_1,\mathcal{R}_1)$ and $(G_2,\mathcal{R}_2)$ are dual: there exist coarse-grainings $C_1, C_2$ such that $C_1(G_1(t))$ and $C_2(G_2(t))$ are equivalent stochastic processes.
\end{quote}

We want:

\begin{itemize}
  \item examples beyond toy models,
  \item invariants preserved under duality,
  \item algorithms for finding duals in rule space.
\end{itemize}

This gives a way to compress ugly universes into prettier ones without losing behavior.

\subsubsection*{30.4  The CΩSMOS Inverse Problem}

We sketched CΩSMOS as "our universe as RMG." Now turn that into a serious inverse problem.

\paragraph{Problem 8: Define a CΩSMOS fitness functional.}

Given a candidate $(G,\mathcal{R})$, define a score $F(G,\mathcal{R})$ that:

\begin{itemize}
  \item heavily rewards matching empirical physics (dimension, symmetries, particle spectra, cosmology),
  \item penalizes weird artifacts (anisotropies, violations of known bounds),
  \item and is computationally estimable from finite-time simulations.
\end{itemize}

Key challenge: making $F$ robust and not insanely fragile to coarse-graining choices.

\paragraph{Problem 9: Search strategies in CΩSMOS-space.}

We’re not going to brute-force MRMW. We need structured search:

\begin{itemize}
  \item parameterized families of rules (graph locality, limited motif types),
  \item gradient-like signals in rule space (how changes in local patterns affect $F$),
  \item priors from existing physics (Lorentz invariance, locality, gauge structure) baked into search.
\end{itemize}

This is not "Evolve cellular automata until they look pretty.’’
This is a disciplined program for finding RMGs that deserve to be CΩSMOS candidates.

\subsubsection*{30.5  Instrumentation: How Do You Actually \emph{Measure} Ruliometry?}

Everything in this book assumes we can probe MRMW as if we had lab instruments.

We don’t. Yet.

\paragraph{Problem 10: Build curvature and dimension estimators.}

Given:

\begin{itemize}
  \item a large RMG evolution (e.g., a CΩMPUTER runtime log),
  \item partial access (e.g., local neighborhoods, sampled histories),
\end{itemize}

design algorithms that estimate:

\begin{itemize}
  \item local dimension (volume growth vs. radius),
  \item curvature proxies (comparisons of geodesic spreads),
  \item bottlenecks and hubs (places where flow concentrates).
\end{itemize}

These should be:

\begin{itemize}
  \item cheap enough to run online,
  \item robust to noise and coarse-graining,
  \item tunable to different scales.
\end{itemize}

\paragraph{Problem 11: Complexity and entropy meters for live systems.}

Given a running system (model, marketplace, codebase), can we:

\begin{itemize}
  \item associate it with a region in MRMW,
  \item track its effective entropy $S(t)$ under a chosen coarse-graining,
  \item estimate its algorithmic complexity (up to practical bounds),
  \item monitor "distance" between its actual trajectory and some safe / desired region?
\end{itemize}

This is the basis for all the "debugging reality" fantasies: you can’t steer what you can’t measure.

\subsubsection*{30.6  Multiversal Machine Theory Proper}

We gave you a handful of multiversal machines in Part~IV.
That’s not a theory, that’s a demo reel.

\paragraph{Problem 12: A calculus of multiversal machines.}

We want:

\begin{itemize}
  \item a formal language for describing machines that operate on MRMW (fork, merge, reweight, prune worlds),
  \item compositionality: you can connect machines and know what the composite does,
  \item expressivity results: which transformations on distributions over worlds are representable.
\end{itemize}

Analogy:

\begin{itemize}
  \item automata theory for string processors,
  \item circuit complexity for Boolean functions,
  \item category theory for compositional structures.
\end{itemize}

Here we want the same, but for operators on measures over MRMW.

\paragraph{Problem 13: Resource theories for multiversal computation.}

Define currencies:

\begin{itemize}
  \item rulial bandwidth (how fast you can move in rule space),
  \item universe budget (how many / how deep worlds you can spin up),
  \item provenance depth (how much history you keep).
\end{itemize}

Then:

\begin{itemize}
  \item characterize which transformations are possible under bounded resources,
  \item prove no-go theorems (impossible steering tasks),
  \item define efficient normal forms for common tasks (simulation, planning, search).
\end{itemize}

This is where "what’s the cost of thinking about this in parallel realities?’’ becomes a formal question.

\subsubsection*{30.7  Intelligence Metrics That Don’t Suck}

We sketched intelligence as efficient, robust steering of MRMW.
Now turn that into metrics that beat "it scored 89.2 on benchmark X."

\paragraph{Problem 14: Define standardized intelligence indices.}

Based on:

\begin{itemize}
  \item reachability (size/shape of reachable good region),
  \item robustness (performance under environment rule perturbations),
  \item model efficiency (compression of local geometry),
  \item multi-scale control (how many levels of coarse-graining the system can reason across).
\end{itemize}

A good index should:

\begin{itemize}
  \item be comparable across architectures,
  \item penalize brittle or overfitted flows,
  \item correlate with qualitative capabilities we care about (transfer, alignment, self-correction).
\end{itemize}

\paragraph{Problem 15: Phase diagrams of intelligence.}

Map out regimes:

\begin{itemize}
  \item below some capacity: flailing, local minima, no reliable generalization;
  \item above thresholds: sudden access to new behaviors (self-modeling, rulial moves, multiverse planning).
\end{itemize}

Formally:

\begin{quote}
Given a family of agents parameterized by capacity (compute, memory, rulial bandwidth), identify critical surfaces in parameter space where qualitative behavior changes.
\end{quote}

This is what everyone hand-waves as "sharp capability jumps.’’
You can do better than vibes.

\subsubsection*{30.8  Ethics and Governance as First-Class Geometry Problems}

We just gave you one chapter on ethics. That was table-setting.

The actual research program is bigger:

\paragraph{Problem 16: Catalog measures over MRMW that correspond to sane ethical stances.}

You need:

\begin{itemize}
  \item concrete candidate measures $\mu$ (how much each world/agent "counts"),
  \item analysis of their failure modes (where they go insane),
  \item approximation schemes for implementing them in real multiversal machines.
\end{itemize}

A good result looks like:

\begin{itemize}
  \item "Measure family $\mathcal{M}_1$ avoids X but fails at Y; $\mathcal{M}_2$ does the opposite; hybrid $\mathcal{M}_3$ dominates both under constraints Z."
\end{itemize}

\paragraph{Problem 17: Formalize rulial constitutions and prove meta-theorems.}

Given predicates $\Phi(\mathcal{R})$ and $\Psi(\mathcal{R}\to\mathcal{R}')$ that encode "allowed universes" and "allowed rule changes," prove things like:

\begin{itemize}
  \item completeness: all safe transformations we care about are permitted;
  \item soundness: no known catastrophic transformation is permitted;
  \item invariance: constitutions themselves are robust under embedding in larger MRMW systems.
\end{itemize}

Also: impossibility theorems.
If certain goals are mutually incompatible (e.g. absolute safety + maximal expressivity), prove it and state the tradeoff precisely.

\subsubsection*{30.9  Building Actual CΩMPUTERs: Experimental Ruliometry}

Theory without experiments is religion.

\paragraph{Problem 18: Minimal CΩMPUTER prototypes.}

Not "hello world’’ nonsense. We want minimal systems that:

\begin{itemize}
  \item represent RMGs natively (graphs of graphs, not forced into arrays),
  \item implement DPO rewriting with provenance (every rewrite event logged as a node in a causal graph),
  \item support basic multiversal operations (fork/bundle/merge) at moderate scale,
  \item expose metrics hooks for curvature/entropy estimates.
\end{itemize}

This is the platform for every empirical question above.

\paragraph{Problem 19: Toy universes as experimental testbeds.}

Design families of RMG rule sets that:

\begin{itemize}
  \item are rich enough to have emergent phenomena (phases, particles, patterns),
  \item are small enough to simulate extensively,
  \item and come with a library of known behaviors.
\end{itemize}

Use them to:

\begin{itemize}
  \item test ruliometric metrics,
  \item evaluate multiversal machines,
  \item prototype ethical constraints in safe sandboxes.
\end{itemize}

Think "Ising model’’ and "Game of Life,’’ but for full-stack MRMW.

\subsubsection*{30.10  Ruliometric Culture: How This Becomes a Field, Not a Meme}

You don’t get a scientific discipline by writing one book.
You get it by building infrastructure:

\paragraph{Problem 20: Shared formalisms and benchmarks.}

We need:

\begin{itemize}
  \item common RMG / DPO specification languages,
  \item shared toy universe suites,
  \item canonical multiversal machine definitions,
  \item standard ruliometric benchmark tasks ("measure curvature here; classify this universe; compare these flows’’).
\end{itemize}

Otherwise everyone reinvents their own private MRMW and nobody’s results talk to each other.

\paragraph{Problem 21: Sociotechnical guardrails.}

Meta, but important:

\begin{itemize}
  \item How do you keep a field that can literally engineer universes from being captured by whichever lab gets richest first?
  \item How do you share CΩMPUTER-like runtimes without handing out multiversal weapons?
  \item What institutional patterns (open consortia, treaty-like agreements, verification infrastructures) make sense here?
\end{itemize}

If "ruliometric science’’ becomes the new nuclear physics and we learn nothing from the last century, that’s on us.

\subsubsection*{30.11  Where to Start (Yes, You, Personally)}

This all sounds huge. It is. That’s the point.
But you don’t have to eat MRMW in one bite.

Some tractable entry points:

\begin{itemize}
  \item \textbf{If you’re a theorist:}
    \begin{itemize}
      \item formalize one specific metric or curvature notion on a restricted MRMW slice (e.g., cellular automata),
      \item prove even one non-trivial theorem relating that to known complexity or dynamical behavior.
    \end{itemize}

  \item \textbf{If you’re a PL / systems person:}
    \begin{itemize}
      \item build a small CΩMPILER prototype: embed a tiny DSL into a graph-rewrite runtime with provenance;
      \item make time-travel debugging or counterfactual execution a first-class feature and see what breaks.
    \end{itemize}

  \item \textbf{If you’re a physicist:}
    \begin{itemize}
      \item pick a toy theory (Ising, scalar field, lattice gauge),
      \item compile it into RMG + DPO cleanly,
      \item and study its MRMW neighborhood: duals, universality classes, curvature.
    \end{itemize}

  \item \textbf{If you’re on the AI / alignment side:}
    \begin{itemize}
      \item recast one alignment problem in MRMW terms (measures, flows, constitutions),
      \item build a tiny multiversal machine that exhibits the failure mode, then one that avoids it.
    \end{itemize}

  \item \textbf{If you’re a philosopher or ethicist:}
    \begin{itemize}
      \item pick a live ethical debate (personhood of sims, population ethics, risk),
      \item express the competing views as different $(v,U,V,\mu)$ choices,
      \item and see where they literally move measure in MRMW.
    \end{itemize}
\end{itemize}

\medskip

The honest answer to "what is ruliometric science right now?’’ is:

\begin{quote}
A sketch of a field that will either quietly die, or become the language we use for physics, computing, AI, and governance in 50 years.
\end{quote}

This book is a bet on the second outcome.

\medskip

If you’ve made it this far, you are now \emph{ruliometric}. You know what MRMW is.
You know what CΩMPUTER is trying to build. You know where the sharp edges are.

The rest is not my problem. It’s yours.

\begin{center}
\textbf{Go bend some universes.}
\end{center}